% main.tex: the ECN 593 applied-project paper (arXiv preprint).
% This folder is a self-contained Overleaf project. Compiler: pdfLaTeX.
% Bibliography: natbib + BibTeX (refs.bib). Locally, from thesis/: latexmk
%
% THIS IS THE TEMPLATE SAMPLE. Every paragraph below is placeholder text that
% exercises one element of ecnthesis.cls. It contains no result numbers.
% Result numbers enter only as macros from numbers.tex (generated; see README).
%
% Look: \documentclass{ecnthesis} is the Kaffeehaus look (default);
%       \documentclass[mono]{ecnthesis} is the monospaced grey look.
\documentclass{ecnthesis}

% ---- generated number macros (scripts/write_numbers.py, copied in by E3).
% Guarded so the template builds before the file exists.
\InputIfFileExists{numbers}{}{\typeout{ecnthesis: numbers.tex not found; number macros not loaded yet}}

% ---- project notation (same as writeup/main.tex)
\newcommand{\NROY}{\mathrm{NROY}}
\newcommand{\hm}{\mathrm{hm}}
\newcommand{\expset}{\mathrm{exp}}
\newcommand{\prem}{\psi}
\newcommand{\interior}{\operatorname{int}}
\newcommand{\Jone}{J_1}
\newcommand{\Jtwo}{J_2}
\newcommand{\dspill}{\delta_{\mathrm{spill}}}
\newcommand{\ddir}{\delta_{\mathrm{dir}}}

% ---- title page
\title{Not-ruled-out images\\of policy contrasts}
\subtitle{A village agent-based model, calibrated by history matching, and the
analysts it can grade}
\shorttitle{Not-ruled-out images of policy contrasts}
\author{Augusto Gonz\'alez-Bonorino}
\affiliation{Arizona State University}
\course{ECN 593 \textperiodcentered\ Applied Project}
\date{Draft, November 2026}
\keywords{partial identification; history matching; interference; agent-based models; calibration}
\jel{C15, C18, C31, C63, D85}
\titlenote{Placeholder title note: repository, replication command and a short
acknowledgement go here. \url{https://github.com/Bonorinoa/ecn593-applied-project}}

\begin{abstract}
Placeholder abstract. This paragraph stands in for a summary of about one
hundred and fifty words: the question, the design, the two analysts, what the
evaluation finds and what a reader should take away. It sits on the title page
between two hairlines, set slightly smaller than the body, and it is followed by
the keywords and the JEL codes. A second sentence is here only so that the block
wraps over several lines and the justification can be judged by eye. A third
sentence closes the paragraph.
\end{abstract}

\begin{document}
\maketitle

%================================================================
\section{Introduction}\label{sec:intro}

Placeholder introduction. A calibrated simulator gives a policy effect for each
parameter point that fits the targets. When many points fit equally well, the
reporting logic of a partially identified set applies \citep{manski2003,tamer2010},
and the object to report is the image of the effect over the points the targets
do not rule out. This sentence also shows a textual citation, as in
\citet{vernon2010}, and a cross-reference to Section~\ref{sec:design}.

A second paragraph shows the first-line indent and the measure. It carries an
inline formula, $\theta(\pi,\pi';c)$, a footnote,\footnote{Footnotes are set in
the body face at footnote size.} and \emph{emphasis}. The text is placeholder;
the paper's own sentences replace it in the writing card.

\begin{statusbox}[Status]
Placeholder disclosure. A short box for the status of a claim, a pre-registration
note, or what the reader should not infer. It uses the accent rule and the wash
fill, and it breaks across pages if it must.
\end{statusbox}

The contributions are listed below, to exercise the list style.
\begin{itemize}
  \item A first contribution, stated in one line.
  \item A second contribution, stated in a sentence long enough to wrap onto a
        second line so the hanging indent can be checked.
  \item A third contribution.
\end{itemize}

%================================================================
\section{The planner's problem}\label{sec:planner}

Placeholder. Under interference, the outcome of person $i$ depends on the whole
assignment vector \citep{hudgens2008}, so the natural estimand is a contrast of
assignment rules,
\begin{equation}\label{eq:theta}
  \theta(\pi,\pi';c)
  \;=\;
  \mathbb{E}_c\!\left[\bar Y_T(\pi)\right]
  \;-\;
  \mathbb{E}_c\!\left[\bar Y_T(\pi')\right],
\end{equation}
where $c$ indexes a complete data-generating process. Equation~\eqref{eq:theta}
is placeholder notation carried over from the current write-up.

\begin{definition}[Not-ruled-out set]\label{def:nroy}
Placeholder. Given an implausibility measure $I(c)$ and a cutoff $\tau$, the
not-ruled-out-yet set is $\NROY = \{c \in C : I(c) \le \tau\}$.
\end{definition}

\begin{assumption}[Exposure mapping]\label{ass:exposure}
Placeholder. Outcomes depend on the assignment only through own treatment and an
exposure summary of the neighbours' treatments \citep{aronow2017}.
\end{assumption}

\begin{remark}
Placeholder. Definition~\ref{def:nroy} is not a posterior, and
Assumption~\ref{ass:exposure} is a modelling choice of one analyst, not a property
of the simulator.
\end{remark}

\subsection{Two analysts}\label{sec:analysts}

Placeholder subsection. It shows the second heading level, which carries its
number in the accent colour. Each analyst reports an interval, and the evaluation
grades the interval against the truth and against the image.

\paragraph{A run-in paragraph heading} Placeholder text after a run-in heading,
for short labelled passages inside a subsection.

%================================================================
\section{Design}\label{sec:design}

Placeholder. The order of operations is fixed before any run, as
Figure~\ref{fig:pipe} shows. The diagram uses the shared TikZ styles of the
class, so it takes the colours of whichever look is active.

\begin{figure}[htb]
\centering
\begin{tikzpicture}[ecn, node distance=7mm and 6mm]
  \node[ecn box, text width=100mm] (freeze) {1~~fix $f$, $C$, $z$, $\tau$ and the estimands};
  \node[ecn box, text width=100mm, below=of freeze] (hm)
    {2~~history matching: $\NROY=\{c:I(c)\le\tau\}$\\[0.15em]
     {\scriptsize not a posterior; an empty set is a finding}};
  \node[ecn box, text width=46mm, below=of hm, xshift=-27mm] (A)
    {\textbf{3~~Leg A: image}\\[0.15em]$\Theta_{\hm}=\{\theta(c): c\in\NROY\}$};
  \node[ecn box, text width=46mm, below=of hm, xshift=27mm] (B)
    {\textbf{4~~Leg B: two analysts}\\[0.15em]an interval at one point $c_0$};
  \node[ecn mark, text width=100mm, below=13mm of hm |- A.south] (ov)
    {overlay: each interval against the truth and against its own image};
  \draw[ecn arrow] (freeze) -- (hm);
  \draw[ecn arrow] (hm.south) -- ++(0,-0.25) -| (A.north);
  \draw[ecn arrow] (hm.south) -- ++(0,-0.25) -| (B.north);
  \draw[ecn arrow] (A.south) -- ++(0,-0.3) -| (ov.north);
  \draw[ecn arrow] (B.south) -- ++(0,-0.3) -| (ov.north);
\end{tikzpicture}
\caption{Order of operations (placeholder diagram). Legs A and B are two reports
and do not merge into one object.}
\label{fig:pipe}
\end{figure}

Placeholder paragraph after the figure. The design section of the paper lists
the exogenous quantities, the targets and the cutoff, each with its source in the
configuration files.

%================================================================
\section{Results}\label{sec:results}

Placeholder. Table~\ref{tab:verdicts} shows the zebra table style: booktabs
rules in the table-rule colour, a light stripe on every second body row, and
column heads in the heading face. Every cell here is a placeholder symbol; in the
paper, each number is a macro from \texttt{numbers.tex}.

\begin{table}[t]
\centering
\caption{Verdicts by analyst and diagnostic point (placeholder).}
\label{tab:verdicts}
\begin{zebratabularx}{\linewidth}{Xccc}
\toprule
\colhead{Analyst and quantity} & \colhead{Point $c^\star$} & \colhead{min-I} & \colhead{median-I} \\
\midrule
Foil 1: interval        & $[\,\cdot\,,\,\cdot\,]$ & $[\,\cdot\,,\,\cdot\,]$ & $[\,\cdot\,,\,\cdot\,]$ \\
Foil 1: contains truth  & ---  & ---  & ---  \\
Foil 1: verdict         & ---  & ---  & ---  \\
Foil 2: interval        & $[\,\cdot\,,\,\cdot\,]$ & $[\,\cdot\,,\,\cdot\,]$ & $[\,\cdot\,,\,\cdot\,]$ \\
Foil 2: contains truth  & ---  & ---  & ---  \\
Foil 2: verdict         & ---  & ---  & ---  \\
Image $\Theta_{\hm}$    & \multicolumn{3}{c}{$[\,\cdot\,,\,\cdot\,]$} \\
\bottomrule
\end{zebratabularx}
\tablenote{Notes: placeholder. Source: \texttt{results/*.json} through
\texttt{numbers.tex}. A dash marks a cell the writing card fills with a macro.}
\end{table}

Figure~\ref{fig:image} shows a figure slot. The command
\verb|\includefigure{image_points.png}| includes \texttt{figures/image\_points.png}
when it exists and draws a labelled placeholder until then, so the paper compiles
before the figures are generated.

\begin{figure}[t]
\centering
\includefigure[width=0.92\linewidth]{image_points.png}
\caption{The image of the policy effect over the sampled points (placeholder
caption). Captions sit below figures and above tables.}
\label{fig:image}
\end{figure}

Placeholder paragraph. A results section would discuss the table and the figure
here, in the order the reader meets them, and point to Appendix~\ref{app:extra}
for robustness material.

%================================================================
\section{Conclusion}\label{sec:conclusion}

Placeholder conclusion. One paragraph restates the question, one the answer,
and one what the procedure is not.

\begin{enumerate}
  \item A first limitation, in one line.
  \item A second limitation, long enough to wrap so the numbered list's hanging
        indent can be checked by eye.
\end{enumerate}

%================================================================
\bibliographystyle{aer}
\bibliography{refs}

%================================================================
\appendix

\section{Additional material}\label{app:extra}

Placeholder appendix. Sections in the appendix are lettered, start on a new page,
and number their equations, tables and figures as A.1, A.2 and so on, as in
\begin{equation}\label{eq:app}
  \Theta_{\hm} \;=\; \bigl[\,\min_{c\in\NROY}\theta(c),\; \max_{c\in\NROY}\theta(c)\,\bigr].
\end{equation}
Equation~\eqref{eq:app} is placeholder notation.

\begin{table}[h]
\centering
\caption{An appendix table (placeholder).}
\label{tab:app}
\begin{zebratabular}{lcc}
\toprule
\colhead{Setting} & \colhead{Value} & \colhead{Source} \\
\midrule
Number of agents        & ---  & config \\
Seeded share            & ---  & config \\
Horizon                 & ---  & config \\
Cutoff $\tau$           & ---  & config \\
\bottomrule
\end{zebratabular}
\end{table}

\begin{remark}
Placeholder. Theorem-like environments keep one counter each throughout the
paper, including the appendix.
\end{remark}

\subsection{An appendix subsection}

Placeholder text for an appendix subsection, to check the heading level and the
running header, which reads ``Appendix A'' on these pages.

\end{document}
